\section{How open the declared open flights are}\label{sec:opening}

Why does the mean squared displacement bend down at long lags?

\paragraph{Opening.} The \emph{opening} of a flight is the distance from its takeoff to its landing, its
first and last fix, over the path flown: 1 for a straight line, near 0 for a flight that comes back to its
takeoff.

\begin{figure}
\centering
\includegraphics[width=\textwidth]{open_declared/opening.pdf}
\caption{Opening of the declared open flights, takeoff to landing over the path flown. Left: every declared
open flight and the flights of the datasets of Figure~\ref{fig:declared}. Right: takeoffs below and above
\MountainM~m. Solid, the threshold \MinOpening; dotted, the stricter \Stricter{} that we do not use.}
\label{fig:opening}
\end{figure}

\begin{table}
\centering\small
\begin{tabular}{@{}l r r r r r r r r@{}}
\toprule
 & \multicolumn{4}{c}{flights} & \multicolumn{4}{c}{median opening} \\
\cmidrule(lr){2-5}\cmidrule(l){6-9}
dataset & all & $\Sl$ & $\G$ & $\A$ & all & $\Sl$ & $\G$ & $\A$ \\
\midrule
every declared open flight & 66\,845 &  &  &  & 0.32 &  &  &  \\
lags up to 1\,000~s & 65\,882 & 52\,187 & 13\,673 & 11\,390 & 0.32 & 0.33 & 0.28 & 0.25 \\
lags up to 4\,150~s & 40\,169 & 31\,694 & 3\,768 & 5\,841 & 0.27 & 0.28 & 0.30 & 0.21 \\
lags up to 11\,450~s & 6\,785 & 5\,775 & 158 & 853 & 0.23 & 0.23 & 0.29 & 0.19 \\
\bottomrule
\end{tabular}

\caption{Flights and median opening of the datasets of Figure~\ref{fig:declared}, by population; a flight of
$\G$ and $\A$ counts in each population where it has a piece long enough.}
\label{tab:opening}
\end{table}

\paragraph{The long lags rest on folded flights.} The declared open flights are far from straight lines: their
median opening is \MedianEvery{} (Figure~\ref{fig:opening}, Table~\ref{tab:opening}). The longer the lag, the
more folded the flights it rests on: the median opening is \MedianThousand{} for the dataset up to 1000~s,
\MedianFixed{} up to \DeclTauFixed~s and \MedianLongest{} up to \DeclTauLongest~s, and the share of flights
with an opening below \MinOpening{} grows from \BelowThousand\% to \BelowLongest\%. A flight that comes back near its takeoff moves away and then back, and its displacement
stops growing as the lag nears its duration: the bend of Section~\ref{sec:declared} is the shape of the
flights. In the longest dataset the flights of $\G$ are more open than those of $\Sl$ and $\A$ (median
\MedianLongestG{} against \MedianLongestS{} and \MedianLongestA), so part of the straighter progress of $\G$
at long lags is their shape.

\paragraph{The threshold.} $H$ has a meaning where the mean squared displacement is a power law, a straight
line in log--log, and we want it on flights that are open. A threshold on the opening drops the declared
open flights that are not. A strict one, \Stricter, keeps only straight flights and raises $H$ by
construction; it also keeps \TakeoffLowStrict\% of the flights that take off below \MountainM~m but only
\TakeoffHighStrict\% of those above, where relief and airspace do not let a flight go straight
(Figure~\ref{fig:opening}, right), so it would select the region as well as the shape. We set the
threshold at \MinOpening: the lowest that drops the flights that come back near their takeoff, it keeps
nearly the same share of the two, \TakeoffLowKeep\% and \TakeoffHighKeep\%. It keeps \Kept{} flights,
\KeptShare\% of the declared open ones.

\paragraph{What we keep.} From here on an \emph{open flight} is a declared open flight with an opening above
\MinOpening.
